This research note is the third of three notes from the views-lab00 project. It documents an evaluation-side response to a problem identified in the autoresearch experiments and discussed in Notes A and B of this triptych: the standard Continuous Ranked Probability Score \citep{gneiting2007strictly} treats observations as known with infinite precision, and this assumption is wrong for data with substantial measurement error.

The problem became visible in autoresearch round 7, where I trained a model with two independent uncertainty heads (one trained from gradient utility, one from heteroscedastic Gaussian negative log-likelihood). The two heads converged to a near-perfect anti-correlation across all three violence targets ($r = -0.85$ to $-0.96$). This was not because the architecture failed to differentiate them — it was because the evaluation metric, CRPS, rewards calibrated sharpness regardless of which uncertainty signal you train against. The model learned the calibration that minimised CRPS, and both heads found their way to the same answer.

The deeper problem: CRPS uses an indicator function $\mathbf{1}\{y \le z\}$ that treats the observation $y$ as a point. For UCDP/GED fatality counts -- and for any dataset where the observation is itself a noisy estimate of some underlying quantity -- this is a fiction. A reported value of 47 fatalities is not the truth; it is a best estimate from media coding, witness reports, and analyst judgment, with known sources of uncertainty documented in the dataset codebook \citep{ucdp_codebook}. The "true" value is a draw from a distribution centred on 47, not the value 47 itself.

This note proposes replacing the indicator function with a smooth observational CDF $F_{obs}(z)$ that captures uncertainty about the observation itself. The work is operationalised as a small Python module \texttt{src/lab\_fuzzy\_crps/} with reference implementations of the metric, three constructors for $F_{obs}$ (including one based on the Reported-Value Inflated Gumbel mixture of \citealt{vesco2026uncertainty}), a numerical consistency verification, and a side-by-side demonstration showing how the ranking of forecasters changes.
